% ---- pg-460 （书页 448） ----

\chapter{判别与分类}
\phantomsection\label{chap:11}

\section{引言}  % 原书 11.1
\phantomsection\label{sec:11.1}

判别与分类是两类多元分析方法, 用于将对象(或观测值)分成不同的集合或者将新的对象(或观测值)分配到事先分好的各组之中. 判别分析本质上具有探索性. 作为一种分割方法, 它通常在因果关系不甚明了的情况下被一次性地用来对所观察到的差别进行调查. 分类方法从下述意义上说较少具有探索性: 它们导出一些明确确定的法则, 可用于分配新的对象. 通常分类比判别要求更多的问题结构.

这样, 判别与分类的直接目标可分别叙述如下:

目标 1: 用图形(三维或更低维的)方法或代数方法描述来自若干个已知集合(总体)的对象(或观测值)的差异性. 我们试图找到这样一些“判别量”, 根据这些判别量的数值能尽可能地将这个集合分离.

目标 2: 将对象(观测值)分为两个或更多个类. 重点在于导出一种能最优地将新的对象指派给所分各类的法则.

我们约定用“判别”一词来指目标 1. 这个术语是费希尔(R. A. Fisher, 见文献[9])在分离问题的首次现代处理中引进的. 然而对这个目的来说更具描述性的词是“\mvbold{分离}”. 我们用“\mvbold{分类}”或“\mvbold{分配}”来指目标 2.

一个能将对象分离的函数有时可发挥分配量的作用; 反之, 一个分配对象的法则可以暗示某种判别方法. 在实践中, 目标 1 和目标 2 往往是重叠的, 这时分离与分配之间的区别就变得很模糊.

\section{两个总体的分离与分类}  % 原书 11.2
\phantomsection\label{sec:11.2}

为说明问题, 让我们列出一些情形. 在这些情形下, 人们感兴趣的是: (1)将两类对象分离, 以及 (2)将一个新对象指派给两个类别之一(或两个类别). 为方便计, 将这两个类别分别记为 $\pi_1$ 和 $\pi_2$. 通常根据比如说 $p$ 个相联系的随机变量 $\mathbf{X}'=[X_1,X_2,\cdots,X_p]$ 的测量值来对这些对象进行分离或分类. $\mathbf{X}$ 的观测值从一个类别到另一个类别会有一定程度的差别.\footnote{若 $\mathbf{X}$ 值对于 $\pi_1$ 和 $\pi_2$ 中的对象来说没有很大差别, 则两个类将不可区分, 因而可将一个新对象不加区分地指派给任一类别.}我们可以将第一类中的全部数值设想为对于 $\pi_1$ 的 $\mathbf{x}$ 值总体; 而将第二类中的全部数值设想为对于 $\pi_2$ 的 $\mathbf{x}$ 值总体. 于是这两个总体可以用概率密度函数 $f_1(\mathbf{x})$ 和 $f_2(\mathbf{x})$ 来描述. 因此我们可以交
